
\section{Numerical Representation and Floating-Point Arithmetic}

Digital computer systems operate with finite storage and finite numerical
representations. Consequently, only a finite subset of the real numbers can be
represented exactly by a computer
\cite{goldberg1991floating,higham2002accuracy}.

Numerical representation in digital systems is based primarily on positional
number systems, particularly the binary system. In general, a number in base
$b$ may be represented by

\begin{equation}
    x =
    \sum_{i=m}^{n} d_i b^i,
\end{equation}

where $d_i$ denotes a digit in the chosen base.

Modern computer systems commonly represent non-integer numerical values using
floating-point formats. In simplified form, a floating-point number may be
written as

\begin{equation}
    x = \pm m \times b^e,
\end{equation}

where $m$ represents the significand, $b$ denotes the numerical base, and $e$
represents the exponent.

The term \textit{significand} is generally more precise than \textit{mantissa}
in the context of modern floating-point arithmetic, although both terms are
frequently encountered in the literature.

Because only a finite number of bits is available for representing these
components, many real numbers cannot be stored exactly. Instead, the computer
stores a nearby representable value
\cite{goldberg1991floating,higham2002accuracy}.

This phenomenon is particularly evident for decimal fractions whose expansions
do not terminate in binary representation. For example, values such as
$0.1_{10}$ cannot be represented exactly using a finite binary expansion.

Floating-point arithmetic therefore introduces rounding errors as a natural
consequence of finite representation. These errors may interact with other
numerical phenomena, including loss of significance, cancellation, underflow,
overflow, and propagation of rounding error
\cite{goldberg1991floating,higham2002accuracy}.

The IEEE 754 standard defines widely adopted floating-point formats and
operations, including rules for binary and decimal arithmetic, rounding,
special values such as infinity and NaN, and exceptional numerical conditions
\cite{ieee7542019}.

In iterative numerical algorithms, finite-precision effects are particularly
important because the result obtained during one iteration frequently becomes
the input to the next iteration. Consequently, rounding errors introduced at
an early stage may propagate throughout the computation.

Depending on the stability and conditioning of the numerical problem, such
errors may remain small, accumulate, or be amplified during successive
operations
\cite{higham2002accuracy}.

The study of floating-point representation is therefore essential for the
interpretation of the experimental results obtained in this work. Differences
between implementations should not be evaluated exclusively in terms of
execution time, but also in terms of numerical behavior and finite-precision
effects.
